% ══════════════════════════════════════════════════════════════════════════════
%  EXPERIMENT 1 — IS YEAST ALIVE?                         (Lesson 1, T2 Week 9)
% ══════════════════════════════════════════════════════════════════════════════
\experimentdiv{Is Yeast Alive?}{Testing for a life process \textbullet{} A control \textbullet{} Measuring a balloon over time}{1}\label{sec:1}
\lessoncodes{VC2S8U01 \textbullet{} VC2S8I01, I02, I04, I06}

\vocabhead
\begin{ctab}
\begin{tabular}{|SW{4.2cm}|W{11.6cm}|}
  \hline\hd & \thead{Word} & \thead{Definition (copy from the slide)} \\ \hline
  & \scaf{\textbf{Life process}} & \scaf{Something every living thing does: movement, respiration, sensitivity, growth, reproduction, excretion, nutrition (MRS GREN).} \\ \hline
  & \textbf{Respiration} & \\ \hline
  & \textbf{Micro-organism} & \\ \hline
  & \textbf{Control} & \\ \hline
  & \textbf{Circumference} & \\ \hline
\end{tabular}
\end{ctab}

\begin{lipanel}
\inquiryli{Write a hypothesis, use a control and record measurements over time.}
\understandli{Use evidence of a life process to decide whether yeast is living.}
\sctitle
\begin{itemize}[leftmargin=1.4em, itemsep=0pt, topsep=1pt]
  \item I can write a hypothesis that names both variables.
  \item I can explain what my control flask shows.
  \item I can use my results to argue that yeast is alive.
\end{itemize}
\end{lipanel}

\begin{skillsfocus}
\focus{QP1}{Write a testable hypothesis.}
\focus{PC2}{Name the variables and the control.}
\focus{PC6}{Measure each balloon every five minutes.}
\focus{PA1}{Record measurements in a THUD table.}
\focus{EV1}{Say whether the results support the hypothesis.}
\end{skillsfocus}

% ─── TECHNIQUE ────────────────────────────────────────────────────────────────
\exppart{Technique spotlight}{Blue}

\begin{technique}{What a control is for}{PC2}
A \textbf{control} is a set-up that is the same as the others except for the one thing you are testing. If the yeast with sugar blows up its balloon and the yeast \emph{without} sugar does not, the sugar made the difference, not the warm water or the flask. Without a control you cannot tell what caused the result.
\end{technique}

\begin{technique}{Measuring a round object}{PC6}
\begin{minipage}[c]{0.62\linewidth}
Wrap a piece of string once around the widest part of the balloon. Pinch the string where it meets its end, then lay it straight along a ruler. Read the length to the nearest 0.5\,cm. This is the \textbf{circumference}. Measure the same way, at the same place, every time.
\end{minipage}\hfill
\begin{minipage}[c]{0.34\linewidth}\centering
\begin{tikzpicture}[scale=0.7]
  \draw[eq] (0,0) .. controls (-1.3,0.6) and (-1.2,2.4) .. (0,2.6) .. controls (1.2,2.4) and (1.3,0.6) .. (0,0);
  \draw[eq] (-0.15,0) -- (0.15,0) -- (0.1,-0.25) -- (-0.1,-0.25) -- cycle;
  \draw[Ink, line width=1.2pt, dashed] (-1.02,1.4) .. controls (-0.5,1.15) and (0.5,1.15) .. (1.02,1.4);
  \eqlab{(1.0,1.4)}{(2.0,2.0)}{string}
\end{tikzpicture}
\end{minipage}
\end{technique}

\begin{workedex}{Writing the hypothesis}
\emph{If sugar is added to the yeast, the circumference of the balloon will increase.}

\smallskip
With a reason: \emph{\dots because the yeast uses the sugar for respiration and gives off carbon dioxide gas.}
\end{workedex}

\begin{wedo}{Writing the hypothesis}
Write a hypothesis for flask C, which has sugar and yeast but \emph{cold} water.
\answerbox{3cm}
\end{wedo}

% ─── PLAN AND SAFETY / METHOD ────────────────────────────────────────────────
\exppart{Plan, safety and method}{Teal}

\textbf{Aim:} To investigate the effect of changing the food and temperature of yeast on the circumference of a balloon over its flask.

\task Complete the variables table. \skr{PC2}
\begin{ctab}
\begin{tabular}{|SW{4.4cm}|W{10.2cm}|}
  \hline\hd & \thead{Variable} & \thead{In this experiment} \\ \hline
  & \scaf{Independent} & \scaf{What is in the flask: sugar or no sugar; warm or cold water} \\ \hline
  & Dependent & \\ \hline
  & Controlled & \\ \hline
  & Controlled & \\ \hline
  & Control flask & \\ \hline
\end{tabular}
\end{ctab}

\needspace{8\baselineskip}
\task Write your hypothesis for flask A. \skr{QP1}
\answerbox{2.5cm}

\subsubsection*{Equipment}
3 conical flasks (250\,mL) labelled A, B and C \textbullet{} 3 balloons \textbullet{} 3 teaspoons of dried yeast \textbullet{} 2 teaspoons of sugar \textbullet{} warm water (about 40\,\textdegree C) and cold water \textbullet{} 100\,mL measuring cylinder \textbullet{} funnel \textbullet{} string and a ruler \textbullet{} stopwatch

\begin{center}
\begin{tikzpicture}[scale=0.72]
  \foreach \x/\l/\c in {0/A/{warm water, sugar, yeast},4.2/B/{warm water, yeast},8.4/C/{cold water, sugar, yeast}}{
    \begin{scope}[shift={(\x,0)}]
      \eqFlask
      \fill[black!12] (-0.12,0.03) -- (0.1,0.7) -- (1.7,0.7) -- (1.72,0.03) -- cycle;
      \draw[eq] (0.55,2.95) .. controls (0.1,3.6) and (0.3,4.5) .. (0.8,4.6) .. controls (1.3,4.5) and (1.5,3.6) .. (1.05,2.95);
      \node[font=\headingfont\bfseries] at (0.8,1.5) {\l};
      \node[eqtext, align=center, text width=3.2cm] at (0.8,-0.55) {\c};
    \end{scope}}
\end{tikzpicture}
\expfigcap{Flask B has no sugar; flask C is cold.}
\end{center}

\needspace{12\baselineskip}
\task Complete the risk assessment. \skr{PC3}
\begin{ctab}
\begin{tabular}{|SW{4.0cm}|W{5.2cm}|W{5.4cm}|}
  \hline\hd & \thead{Risk (what could go wrong)} & \thead{Prevention} & \thead{Responding (if it happens)} \\ \hline
  & \scaf{A glass flask breaks} & \scaf{Keep flasks in the middle of the bench; carry one at a time} & \scaf{Do not pick it up; tell the teacher, who sweeps it up} \\ \hline
  & Warm water spills & & \\ \hline
  & Someone is allergic to latex balloons & & \\ \hline
\end{tabular}
\end{ctab}

\subsubsection*{Method}
\begin{enumerate}[itemsep=2pt]
  \item Using the funnel, put 1 teaspoon of yeast into each flask, and 1 teaspoon of sugar into flasks A and C.
  \item Add 100\,mL of warm water to flasks A and B, and 100\,mL of cold water to flask C. Swirl each gently.
  \item Stretch a balloon over the neck of each flask straight away, and start the stopwatch.
  \item Measure the circumference of each balloon at 0, 5, 10, 15 and 20 minutes. Record each in Table 1.1.
  \item Pour the contents down the sink with plenty of water, and rinse the flasks.
\end{enumerate}

% ─── RESULTS ─────────────────────────────────────────────────────────────────
\exppart{Results}{Amber}

\subsubsection*{Table 1.1: Circumference of each balloon over time \skr{PA1}}
\begin{ctab}
\begin{tabular}{|SW{2.6cm}|C{2.2cm}|C{2.2cm}|C{2.2cm}|C{2.2cm}|C{2.2cm}|}
  \hline\hd & \thead{Flask} & \thead{0 min (cm)} & \thead{5 min (cm)} & \thead{10 min (cm)} & \thead{15 min (cm)} & \thead{20 min (cm)} \\ \hline
  & \scaf{Example} & \scaf{6.0} & \scaf{8.5} & \scaf{12.0} & \scaf{15.5} & \scaf{18.0} \\ \hline
  & A & & & & & \\ \hline
  & B & & & & & \\ \hline
  & C & & & & & \\ \hline
\end{tabular}
\end{ctab}

\needspace{12\baselineskip}
\task Describe what you saw in each flask, apart from the balloon. \skr{PC6}
\answerbox{4cm}

% ─── ANALYSIS ────────────────────────────────────────────────────────────────
\exppart{Analysis and evaluation}{Coral}

\task Describe the trend in flask A. Use numbers from Table 1.1. \skr{PA5}
\answerbox{3.5cm}

\needspace{10\baselineskip}
\task Flask B is a control. What does it tell you about the result in flask A?
\answerbox{4cm}

\needspace{10\baselineskip}
\task Explain the result in flask C.
\answerbox{3.5cm}

\needspace{12\baselineskip}
\task Which life processes (MRS GREN) did your experiment give evidence for? For each, say what the evidence was.
\answerbox{5cm}

\needspace{12\baselineskip}
\task Write a conclusion: state the trend, back it with numbers from your table, and say whether your hypothesis was supported. \skr{EV1}

\scaf{The yeast with sugar in warm water made the balloon grow from \dots\,cm to \dots\,cm. This supports / does not support the hypothesis.}
\answerbox{5cm}

\begin{yr8box}
\begin{keepext}
\exttask A student boils yeast for five minutes, cools it, and repeats flask A with it. Predict what will happen to the balloon, and explain what this would show about dead and living things.
\answerbox{5cm}
\end{keepext}
\end{yr8box}

% ─── SKILL CHECK ─────────────────────────────────────────────────────────────
\exppart{Skill check}{NavyBlue}

\begin{skillcheck}
\scrow{QP1}{I can write a hypothesis that names both variables.}
\scrow{PC2}{I can explain what a control is for.}
\scrow{PC6}{I measured each balloon the same way every five minutes.}
\scrow{PA1}{My table has a title, headings, units in the headings and matching decimals.}
\scrow{EV1}{I can say whether my results supported the hypothesis.}
\end{skillcheck}

\task Which skill will you focus on in Experiment 2? Explain what you will do differently.
\answerbox{5cm}
